\chapter*{General Introduction}
\addcontentsline{toc}{chapter}{General Introduction}
\markboth{General Introduction}{General Introduction}

Legal services have historically been expensive, complex, and difficult to access,
particularly for individuals and small businesses navigating an increasingly
regulated environment. The rapid development of \ac{AI} is beginning to change
this: language models can now engage in open-ended dialogue about legal matters,
interpret regulations, and assist users in understanding their rights and
obligations. For a LegalTech company like E-Tafakna, which already offers a
comprehensive platform for contract management and legal document processing
across the \ac{MENA} region, this evolution opens the door to a qualitative
improvement in the intelligence it can offer to its users.

Two obstacles, however, stand between the current state of the art and a
genuinely useful Tunisian legal assistant. First, general-purpose language models
were not trained on Tunisian legislation: they lack the specific vocabulary,
article structure, and regulatory context of local law, and their responses on
Tunisian legal matters are correspondingly unreliable. Second, the legal
sensitivity of the data makes external cloud-based inference unacceptable.
E-Tafakna's planned deployment at \ac{ATI} requires a system that runs
entirely on local infrastructure, with no data leaving the controlled environment
during operation.

This graduation project addresses both obstacles simultaneously. We design,
build, and evaluate a domain-specific legal \ac{AI} assistant that combines
a language model fine-tuned on Tunisian legal texts with a
\ac{RAG} pipeline that retrieves relevant legal articles at inference time,
all deployed on-premises with no external dependencies at query time.

This report is organized in four chapters. Chapter~I situates the project within
its institutional and methodological context: it presents E-Tafakna, analyses the
limitations of its current \ac{AI} agents, defines the problem and the proposed
solution, and justifies the adoption of \ac{DSR} as the guiding methodology.
Chapter~II reviews the scientific and technical literature relevant to the
project, covering \ac{RAG} architectures, parameter-efficient fine-tuning, and
on-premises deployment of large language models. Chapter~III describes the
complete design and development of the artifact: the seven-phase offline
preparation pipeline, the real-time query architecture, the fine-tuning of two
candidate models on a curated Tunisian legal dataset, and their quantized local
deployment. Chapter~IV closes the \ac{DSR} cycle by demonstrating the system in
operation and evaluating it against the three objectives defined in Chapter~III,
covering training outcomes, retrieval quality, expert assessment, and on-premises
compliance.
